\begin{afterword}
\summaryhead

The post opens with a conversation. The futurist Greg Stock predicted pills that would simulate the joy of scientific discovery and be more fun than the real thing; the author replied that this was a reason never to take them. Against ethicists who argue as if all desires reduce to the desire for happiness, with Sam Harris as the example, the post asks whether we should care about things apart from the happiness they bring.

Doing what makes us happy does not make happiness the only reason: that would leave others' happiness a means to a warm glow, and an act can have several consequences we value. A reduction would need happiness to be the only thing that counts in every decision. The author would not value art that no one ever sees, would not enter a holodeck even with a pill to forget it, and values freedom from unseen manipulation. The conclusion: the author has many terminal values, none reducible to happiness or to subjective states.

\responsehead

The post states, clearly and briefly, the standard case against hedonism about value, and its conclusion is now the majority view. In the 2020 PhilPapers survey 76.9 per cent of philosophers would not enter an experience machine, and 12.7 per cent leaned toward hedonism about well-being. The cases are older than the post. G. E. Moore judged a life of enjoyment among ``absolutely unreal'' objects ``greatly inferior in value'' in 1903, and held that an emotion and its object have little value apart. Robert Nozick's experience machine (1974) is the holodeck with the forgetting pill. The post does not mention either, and does not claim the cases as new; it credits Sober and Wilson for one point.

The arguments in the middle do not reach the question the post poses. That question is evaluative: whether we should care about more than happiness. The first argument shows that caring about others is hard to explain if one's own happiness is the only reason. That tells against egoism, not against the view the post set out to examine, which counts others' happiness as well. The second, that a result of an action need not be its aim, is Butler's old objection to psychological egoism. The third sets a bar for the claim that happiness is the only thing that counts in our decisions. All three concern motives. An evaluative hedonist can accept them and still hold that only happiness is worth caring about.

What does address the question is the cases, and there the evidence is the author's intuition, offered as such (``my moral intuition appears to require''). The conclusion is stated about ``my values.'' What the \textsc{Stanford Encyclopedia} calls the hedonists' strongest reply is that such intuitions are unreliable guides to value, and the post does not consider it. Later experiments give it some support: people told they were already in an experience machine often chose to stay. On unseen beauty, the post's answer is Sidgwick's; Moore answered the other way.

Last, the step from ``not strictly reducible to happiness'' to values ``none of them strictly reducible to anything else'' rests on a linked earlier post, not on these cases. If the intuitions are accepted, the cases count against hedonism. They fit desire theories of well-being as well as the hybrid view the post's ``both'' describes.

In short: a clear statement of the standard case against hedonism, made earlier by Moore and Nozick; its arguments about motives do not reach the evaluative question it poses, which it answers with the author's own intuitions.
\end{afterword}
